% A.77 candidate. Notes only; no canonical promotion before full audits.
\subsection{Closures with at most two commutator elements}

\begin{statement}{A.77}{Open}{Candidate: both cusp assertions when the geometric commutator has order at most two}
Let \(Y\to E\setminus\{O\}\) be an actual connected finite
\'{e}tale cover whose finite geometric Galois closure group
\(G\) satisfies \(|G'|\le2\). Then both assertions of
Task A.5 hold for \(Y\), at every coefficient prime,
for every original induced cusp packet and every permitted
open derived subgroup. Every actual intermediate is included.

For a nonabelian two-group closure \(G\), its proper
Galois source \(C\) has potential good Jacobian reduction
of special two-rank zero at two. Its cusp inertia is
\(G'\), its genus is \(1+|G|/4\), and its cusp
number is \(|G|/2\). If \(e\) is the exponent of
\(G/G'\), every cusp difference is killed by \(2e\).
The rank-zero and torsion assertions here concern the
two-group source; the odd abelian extension below needs
only the actual rational-path preservation.

This includes, for example, the actual order-thirty-two
closure
\[
 \langle x,y,z\mid x^4=y^4=z^2=1,\quad
             [x,y]=z,\quad [x,z]=[y,z]=1\rangle.       \tag{A.77.1}
\]
Its full proper source has genus nine and sixteen cusps.
Its largest cyclic subgroup has order four, so it is
outside the cyclic-index-at-most-two class of A.76.
The universal question A.45 remains open.
\end{statement}

\begin{proof}[Candidate proof, pending complete independent audits]
An actual geometric Galois closure is a quotient of the
rank-two free profinite group of the punctured elliptic
curve. Thus \(G\) is generated by the actual images
of two peripheral-handle generators. Pass to a finite
number-field extension defining its cover, deck group,
quotients and marked points. All later comparison maps
are actual morphisms of these curves or Jacobians.

If \(G\) is abelian, the cover compactifies to an
elliptic isogeny and A.40 applies. Suppose \(G'\simeq C_2\).
Conjugation on this group is trivial, since
\(\operatorname{Aut}(C_2)=1\). Hence \(G'\) is central,
\(G\) has nilpotency class two, and the finite group
is the direct product of its Sylow subgroups. Its odd
Sylow subgroups are abelian, since their commutators
lie in both an odd group and \(G'\). Write
\[
                            G=G_2\times B,
                                                        \tag{A.77.2}
\]
where \(B\) is odd abelian. Projection of the actual
two generators generates \(G_2\). We first prove the
stronger geometric statements for this two-group and
all its intermediates, then retain the odd factor by
an actual elliptic-isogeny pullback.

\paragraph{Every faithful irreducible quotient has a cyclic subgroup of index two.}
Let \(T\) be a two-generated finite two-group with
\(T'\simeq C_2\), and let \(\rho\) be an irreducible representation over
an algebraic closure of an auxiliary \(\mathbf Q_\ell\),
with \(\ell\ne2\). Put \(N_\rho=\ker\rho\)
and \(Q=T/N_\rho\). This is an actual finite group
quotient, with a faithful irreducible representation.
If it is abelian, it is cyclic, because a faithful
one-dimensional finite scalar image in characteristic zero is cyclic.
Suppose \(Q\) nonabelian. Then \(Q'\) is still
of order two and central. Schur's lemma embeds its
centre \(Z\) in that scalar group, so \(Z\)
is cyclic of order \(2^c\), with \(c\ge1\).

Every square in \(Q\) is central, since
\([a^2,b]=[a,b]^2=1\). The two actual generators
\(u,v\) therefore give an elementary abelian quotient
\(Q/Z\) generated by at most two elements. It cannot
be cyclic, because a group cyclic modulo its centre
is abelian. Consequently
\[
                Q/Z\simeq C_2^2,\qquad |Q|=4|Z|.
                                                        \tag{A.77.3}
\]
Let \(\sigma=[u,v]\) be its nontrivial commutator.
The class-two word collection is exact: every word
is \(u^a v^b\sigma^d\). It commutes with \(u\)
and \(v\) exactly when \(a,b\) are even. Thus
\[
                         Z=\langle u^2,v^2,\sigma\rangle.
                                                        \tag{A.77.4}
\]
If \(c\ge2\), \(\sigma\) belongs to the square
subgroup \(Z^2\). If both \(u^2,v^2\) also belonged
to \(Z^2\), (A.77.4) would give \(Z\subset Z^2\),
which is impossible. One of those squares therefore
generates the cyclic \(Z\), and its square root in
\(Q\) has order \(2|Z|\). It generates a cyclic
subgroup of index two. If \(c=1\), the same argument
works unless \(u^2=v^2=1\); in that case
\((uv)^2=\sigma\), and \(uv\) has order four and
index two in the order-eight group \(Q\).

This proves the claimed faithful-quotient property
without a classification of representations or an
assumption that the original centre is cyclic.
A.76's elementary classification identifies every
nonabelian \(Q\) with one of the already proved
cyclic-index-two families. In this class-two situation
only the relevant small dihedral/quaternion or modular
cases can occur, but that further restriction is unnecessary.

\paragraph{The quotient Jacobians are actual and jointly surject.}
Let \(C\to E\) be the proper Galois source for \(T=G_2\).
For each irreducible representation over that splitting field choose its
kernel \(N\), with repeated kernels used only once.
There are finitely many such normal subgroups. Form the
actual quotient curve and quotient map
\[
                    q_N:C\longrightarrow C_N=C/N.
                                                        \tag{A.77.5}
\]
The map \(C_N\to E\) is the actual Galois cover
with group \(T/N\), not merely a curve with an
abstractly matching representation. If that group is
abelian its proper source is an elliptic isogeny cover.
Otherwise its cyclic-index-two group falls under A.70,
A.72, A.73 or A.75. Those proofs give potential good
Jacobian reduction of special two-rank zero at two.
Thus every \(J_{C_N}\) has this property.

Consider the sum of the actual Jacobian pullbacks,
\[
       q^*:\prod_N J_{C_N}\longrightarrow J_C,
                  \qquad (a_N)_N\longmapsto\sum_Nq_N^*a_N.
                                                        \tag{A.77.6}
\]
It is surjective. To verify this, use rational Tate
homology at any auxiliary prime different from two.
Actual pullback and norm satisfy
\[
                 q_N^*q_{N*}=\sum_{h\in N}h,
          \qquad q_{N*}q_N^*=|N|.                    \tag{A.77.7}
\]
These identities already hold on divisors: the pullback of
 the image of a point counts each point in its orbit with
 its ramification index, exactly the stabilizer multiplicity
 in the sum over \(N\); the degree of every fibre is \(\lvert N\rvert\).
 They therefore descend to the actual Picard pullback and norm.
 Consequently the image of \(q_N^*\) is precisely
the \(N\)-invariant subspace of \(V_\ell J_C\).
After a splitting scalar extension, this finite-group
representation is semisimple. Every irreducible summand
is fixed by its own kernel, and that kernel is one of
the \(N\)'s in (A.77.5). The kernels were chosen in exactly this auxiliary-adic
splitting field. Their invariant images therefore span
all of \(V_\ell J_C\), by descent of a spanning
identity after scalar extension.

The abelian image of (A.77.6) has full rational Tate
homology and hence full dimension, proving surjectivity.
Its kernel may have positive dimension: we do \emph{not}
claim that the whole product is isogenous to \(J_C\).
Poincare complete reducibility instead makes \(J_C\)
an isogeny quotient or factor of that product. Potential
good reduction descends to such a quotient, and after
common good reduction the special two-rank is isogeny
invariant and additive. Thus \(J_C\) has potential
good reduction of special two-rank zero.
The quotient-good-reduction input is precisely the
one used in A.40 and A.68--A.75.
\cite[Corollary 2 to Theorem 1]{fghSerreTate1968}

\paragraph{Actual cusp inertia and functions on the two-group source.}
Since \(T\) is class two and generated by the actual
handle images \(a,b\), its commutator subgroup is
generated by \([a,b]\). It is nontrivial of order
two. The actual peripheral relation therefore gives
\[
 I=T'=\langle\sigma\rangle,
 \qquad |S_C|=|T|/2,
 \qquad g(C)=1+|T|/4.                                 \tag{A.77.8}
\]
Here \(I\) is central. The actual quotient
\(E'=C/I\to E\) is an unramified elliptic isogeny
whose kernel is \(T/I\). Let \(e\) be the exponent
of this two-group. For kernel points \(Q_i,Q_j\),
the divisor \(e(Q_i-Q_j)\) is principal on \(E'\).
The degree-two map \(h:C\to E'\) has one cusp
\(P_i\) above each \(Q_i\), with \(h^*Q_i=2P_i\).
Thus
\[
                        2e(P_i-P_j)\text{ is principal}.
                                                        \tag{A.77.9}
\]

For any proper subgroup \(H<T\), the cusp count of
the actual quotient \(D=C/H\) is \([T:HI]\).
It is at least two, because the commutator subgroup
\(I\) lies in the Frattini subgroup of the finite
two-group. Explicitly every maximal subgroup has index
two and an abelian quotient, so contains \(I\);
if \(HI=T\), a maximal subgroup containing the
proper \(H\) would give a contradiction. The genus
of \(D\) is positive because it maps finitely to
\(E\). Its Jacobian is an actual norm quotient of
\(J_C\), and hence has the same potential-good,
special-rank-zero properties. Norms of (A.77.9) kill
every cusp difference on \(D\) by \(2e\), since
pushforward of a point over an algebraically closed
field has coefficient one, independently of ramification.

A.41's torsion-divisor criterion now gives a rational
path to a distinct companion for every cusp of every
such \(D\), including the marked endpoint components.
A.42 gives all full fixed-inertia conclusions. The
quotient \(H=T\) is the original curve and is already
covered by A.40. A.59 gives all normalizer assertions
for these nilpotent closures.

\paragraph{Retaining the odd abelian factor by the actual path construction.}
Return to (A.77.2). Every covering subgroup splits as
\(H_2\times H_B\), by the coprime-power argument of
A.71. Its odd abelian part is an actual elliptic
isogeny \(\eta:E''\to E\). Let \(Z\) be the
proper \(H_2\)-intermediate of the two-group source.
The two actual function fields have coprime degrees,
with the isogeny field Galois. Their compositum is
therefore the geometrically connected actual pullback
\(Y=Z\times_EE''\), with removed points exactly
the pairs over \(O\).

If \(Z\) has genus one, A.40 applies. Otherwise the
stronger geometric input just proved supplies its
compact-type marked model at two and rational cusp
companion paths. Base the Abel map at a cusp sent
to \(O\), extend it over that model, and compose
with the actual extended homomorphism \(J_Z\to E\).
This gives the morphism of models used to pull back
\(\eta\), including all nodes and inserted rational
components.
\cite[Theorem 1.1 and Proposition 3.2]{fghDupuyRabinoff}
The good odd-degree isogeny is finite \'{e}tale, so
its pullback is a semistable model of the actual
connected source. Separately, the geometric closure
has nilpotent group \(G_2\times B\), good elliptic
base and sole branch section \(O\). Lehr--Matignon
Theorem 6.13, followed by the actual Jacobian norm,
proves that this pullback has tree dual graph.
\cite[Theorem 6.13]{fghLehrMatignon2005}

A connected rational cusp-joining subtree on \(Z\)
maps constantly to \(O\). Over it the \'{e}tale
isogeny pullback is a disjoint union of literal copies
of that subtree. The copy selected by a given marked
point also contains a distinct companion, indexed by
the same kernel point of \(\eta\). The generic
points are distinct because their \(Z\)-coordinates
are distinct. Stabilization cannot insert positive
genus into their rational path. A.42 therefore proves
all fixed-inertia assertions for \(Y\); A.59 proves
all normalizer assertions. This is the actual A.71/A.76
preservation proof with its geometric hypotheses checked,
not a transfer of fixed elements through a quotient.

All finite field extensions concern algebraic data and
reduction. Geometric paths are compared up to inner
conjugacy, and the established A.41--A.42 restriction
argument retains the original packets and open derived
subgroups.

\paragraph{The explicit order-thirty-two illustration.}
The group (A.77.1) has normal forms
\(x^i y^j z^k\), with \(i,j\bmod4\), \(k\bmod2\),
and multiplication
\[
 (i,j,k)(i',j',k')=(i+i',j+j',k+k'+ji')
                                                        \tag{A.77.10}
\]
with the stated moduli. It has centre
\(\langle x^2,y^2,z\rangle\simeq C_2^3\), commutator
\(\langle z\rangle\), and exponent four. Thus its
largest cyclic subgroup has index eight. The four
actual quotients obtained by imposing
\(x^2=z^\epsilon\), \(y^2=z^\delta\), with
\(\epsilon,\delta\in\{0,1\}\), are three dihedral
order-eight groups and one quaternion order-eight group.
Their four degree-two irreducible characters, together
with the sixteen linear characters of \(C_4^2\),
exhaust the representation degrees. This illustrates
the actual quotient-Jacobian spanning mechanism without
replacing the general argument by a finite example.

The exact group, quotient and character checks are
retained in \path{notes/fgh_A77_certificate.py}.
\end{proof}
